% Figure: orbit radius of a charged particle in a uniform field against its
% momentum, r = p/(qB), for three field strengths. The radius is linear in
% momentum at fixed field, so a magnetic sector sorts a beam of fixed
% momentum by charge-to-mass ratio along a position the line predicts.
\begin{figure}[htbp]
\centering
\begin{tikzpicture}[>=Latex]
\begin{axis}[
  width=11.5cm, height=7cm,
  xlabel={momentum $p = mv$ ($10^{-21}\,\si{\kilogram\meter\per\second}$)},
  ylabel={orbit radius $r$ ($\si{\centi\meter}$)},
  xmin=0, xmax=10,
  ymin=0, ymax=25,
  legend pos=north west,
  legend cell align=left,
  font=\scriptsize,
  grid=both,
  grid style={enggray!20},
]
% r = p/(qB). With p in units of 1e-21 kg m/s and q = 1.6e-19 C:
% r (m) = p_val*1e-21 / (1.6e-19 * B). r (cm) = 100 * that.
% coefficient per unit p_val: 100*1e-21/(1.6e-19*B) = 6.25e-1 / B (cm per unit)
\addplot[engblue, line width=1.2pt, domain=0:10, samples=2]
  {0.625/0.2 * x};
\addlegendentry{$B = 0.2\,\si{\tesla}$}
\addplot[engred, line width=1.2pt, domain=0:10, samples=2]
  {0.625/0.5 * x};
\addlegendentry{$B = 0.5\,\si{\tesla}$}
\addplot[engamber, line width=1.2pt, domain=0:10, samples=2]
  {0.625/1.0 * x};
\addlegendentry{$B = 1.0\,\si{\tesla}$}
% mark the worked example: 100 u ion at 1 kV in 0.5 T, p ~ 7.3e-21, r ~ 9.1 cm
\addplot[enggray, only marks, mark=*, mark size=2pt]
  coordinates {(7.3,9.1)};
\node[enggray, anchor=south east, font=\scriptsize] at (axis cs:7.2,9.4)
  {$100\,\text{u}$ ion, $1\,\si{\kilo\volt}$};
\end{axis}
\end{tikzpicture}
\caption[Cyclotron orbit radius against momentum]{Orbit radius of a
charged particle in a uniform magnetic field against its momentum, from
$r = p/(qB)$, for three field strengths. The radius is proportional to
momentum at fixed field, so a beam entering a uniform-field region spreads
along the radius axis by momentum, and since two singly charged ions
accelerated through the same voltage carry momenta in the ratio of the
square roots of their masses, the field separates them by mass. The marked
point is the singly ionised $100\,\text{u}$ molecule of the worked example,
accelerated through $1\,\si{\kilo\volt}$ and orbiting at about nine
centimetres in a half-tesla field. Read the other way, the line is the
design constraint of every storage ring: to hold the orbit radius fixed as
the beam is accelerated, the field must rise in step with the momentum,
which is why a synchrotron ramps its bending magnets in lockstep with the
beam energy.}
\label{fig:vol04ch08:cyclotron-orbit}
\end{figure}
